%%%%%%%%%%%%%%%%%%%%%%%%%%%%%%%%%%%%%%%%%%%%%%%%%%%%%%%%%%%%%%%%%%%%%%%%%%%
% Lehreinheit 10: Datenanalyse und Interpretation
% Quelle: WuZ-Buch, Kap. 15 (LE 10), Abschnitt 10d (Vorhersage-Trick,
%         p-Hacking, Tafelformel), Kap. 5b (Newcomb, Horwich),
%         Kap. 12f (explorative Subgruppenanalyse, nur Zahlen aus dem Buch)
% Staunmoment: Die unveraenderliche Summe (selbstarbeitende Variante,
%              Vorschlag) + Vorhersage-Trick
% Kernbegriff: Determinismus
%%%%%%%%%%%%%%%%%%%%%%%%%%%%%%%%%%%%%%%%%%%%%%%%%%%%%%%%%%%%%%%%%%%%%%%%%%%
\nsl{In LE~8 und LE~9 haben die Studierenden Daten erhoben und
beschrieben. Diese Einheit fragt, was man aus ihnen schließen darf. Sie
beginnt mit einer Vorhersage, die versiegelt im Raum liegt, und mit vier
Zahlen, die Freiwillige frei wählen. Die Summe trifft die Vorhersage. Die
Wahl war frei, das Ergebnis trotzdem festgelegt: durch die Struktur der
Rechnung. Dieselbe Struktur steckt in vielen Datenauswertungen, die
scheinbar offen beginnen und doch nur zu einem Ergebnis führen können.
Die Einheit unterscheidet deshalb sauber zwischen Daten, Befund und
Schlussfolgerung, zeigt mit dem Vorhersage-Trick aus dem Buch, wie
p-Hacking garantierte Muster erzeugt, und stellt die Gegenmittel vor:
Gegenhypothesen, Triangulation und Präregistrierung.
\newpage}
\os{\begin{description}
\item[{\rot\bf Heute}]~\\[-\lblskip]
\bi
  \ite eine Summe, die vorher feststand
  \ite Daten, Befund, Schlussfolgerung
  \ite Over-Interpretation, p-Hacking, HARKing
  \ite der Garten der sich verzweigenden Wege
  \ite Gegenmittel: Triangulation und Präregistrierung
\ei
\end{description}
\newpage}

\begin{vorbereitung}[Die unveränderliche Summe und der Vorhersage-Trick]
\textbf{Hinweis.} Kap.~15 des Buchs nennt den Effekt \enquote{Die
unveränderliche Summe (Datum-Force)}, beschreibt aber keine Methode. Die
folgende selbstarbeitende Variante ist ein \emph{Vorschlag} für die
Umsetzung. Sie beruht auf reiner Arithmetik und braucht keine
Fingerfertigkeit.

\textbf{Das Prinzip.} Eine Person nennt (1)~ihr Geburtsjahr, (2)~das
Alter, das sie im laufenden Kalenderjahr erreicht, (3)~das Jahr eines
persönlich wichtigen Ereignisses und (4)~die Anzahl der Jahre, die seitdem
bis zum laufenden Jahr vergangen sind. Geburtsjahr plus erreichtes Alter
ergibt das laufende Jahr $J$, Ereignisjahr plus vergangene Jahre ebenfalls.
Die Summe aller vier Zahlen ist deshalb immer $2J$.

\textbf{Eine Woche vorher}
\begin{checkliste}
  \item \textbf{Termin prüfen!} LE~10 liegt im Dezember oder Januar. Die
    Vorhersage hängt am Kalenderjahr des Vorlesungstags: im Dezember~2026
    lautet sie 4052, im Januar~2027 lautet sie 4054. Wird der Termin über
    den Jahreswechsel verschoben, Vorhersage neu schreiben.
  \item Vorhersage auf eine Karte schreiben, in einen Umschlag stecken,
    zukleben und unterschreiben lassen (z.\,B.\ von einer Kollegin). Der
    Umschlag liegt ab Beginn der Stunde sichtbar am Pult oder klebt an der
    Tafel.
  \item Sprechtext üben, vor allem die Fragen~2 und~4 (siehe
    Fallstricke). Zwei Probeläufe mit Personen, die den Effekt nicht
    kennen; dabei auf Missverständnisse achten.
  \item Für den Vorhersage-Trick (Buch, Abschnitt~10d) die geheime
    Botschaft der Abweichungen festlegen, z.\,B.\ $r=(3,1,4,1)$ für die
    ersten Ziffern von $\pi$. Die letzte Zahl der Runde wird geforct. Das
    Buch nennt dafür Force-Bag, Force-Deck, einen Rechen-Force oder
    Equivoque \parencite{pailhes2020}; Wahl und Einstudierung des Forces
    nach der Standardliteratur, hier bewusst nicht beschrieben. Die
    Master-Vorhersage $p_n = F + r_n$ vorab auf eine Karte schreiben und
    versiegeln.
  \item Plan~B ohne Force aus dem Buch bereithalten: der 1089-Trick oder
    Kaprekars Konstante 6174 (Buch, Abschnitt~10d).
\end{checkliste}

\textbf{Fallstricke der unveränderlichen Summe}
\begin{checkliste}
  \item \textbf{Alter:} Es muss das im laufenden Jahr \emph{erreichte}
    Alter sein, egal ob der Geburtstag schon war. Wer noch nicht
    Geburtstag hatte und das aktuelle Alter nennt, liegt um eins zu
    niedrig. Formulierung: \enquote{Wie alt werden Sie dieses Jahr an
    Ihrem Geburtstag -- oder sind es schon geworden?}
  \item \textbf{Vergangene Jahre:} gezählt wird nach Jahreszahlen, nicht
    nach Jahrestagen. Formulierung: \enquote{Wie viele Silvester liegen
    seit diesem Ereignis hinter Ihnen?} Ein Ereignis im laufenden Jahr
    ergibt null.
  \item \textbf{Jahreswechsel:} Im Januar denken viele noch im alten Jahr.
    Dann \enquote{in diesem Jahr} durch \enquote{im neuen Jahr} ersetzen
    und bei der Silvester-Frage \enquote{das letzte mitgezählt} ergänzen.
    Im Dezember besteht die Gefahr nicht.
  \item \textbf{Paare gehören zusammen:} Geburtsjahr und Alter müssen von
    derselben Person stammen, Ereignisjahr und vergangene Jahre vom
    selben Ereignis. Mehrere Personen sind möglich, wenn jede ein
    vollständiges Paar liefert.
  \item \textbf{Rechnen:} eine dritte Person addiert mit dem
    Taschenrechner, eine vierte kontrolliert. Kopfrechnen vor Publikum
    erzeugt Fehler.
\end{checkliste}

\textbf{Handout und Übungsblatt}
\begin{checkliste}
  \item Handout \enquote{Datenanalyse und Interpretation} sowie
    Übungsblatt LE~10 Teil~A und Teil~B je einmal pro Person kopieren
    (Druckvorlage \texttt{handouts.pdf}, Seiten zu LE~10). Teil~A enthält
    die vollständigen Rohdaten der fiktiven Studie (Aufgabe~1) und liegt
    zu Beginn aus. Teil~B enthält die Tabelle zur Formel
    $1-(1-\alpha)^m$ (Aufgabe~2), die Reflexion und das Tagebuch; er wird
    erst nach der Tafel-Auflösung des Vorhersage-Tricks ausgeteilt, weil
    Aufgabe~2 diese Auflösung und das Tagebuch den Force der ersten
    Vorführung vorwegnähmen.
  \item Taschenrechner zulassen; wer einen Laptop hat, darf die Rohdaten
    in eine Tabellenkalkulation übernehmen.
\end{checkliste}

\textbf{Am Tag selbst}
\begin{checkliste}
  \item Umschlag mit der Vorhersage vor Beginn sichtbar aufhängen; nicht
    kommentieren.
  \item Tafel: rechts oben Platz für die vier Zahlen und die Summe, links
    Platz für die Tabelle des Vorhersage-Tricks (Runde, $a_i$, $p_i$,
    $p_i-a_i$) und die Formel.
  \item Für den Vorhersage-Trick: fünf Karten, einen Umschlag, die
    versiegelte Master-Karte, das Force-Requisit.
\end{checkliste}

\textbf{Risiken und Plan B.} Stimmt die Summe nicht, zuerst die beiden
Fallstricke prüfen (Alter, Silvester), und zwar offen: Das ist dann
selbst ein Lehrstück über Messfehler und Operationalisierung. Der Effekt
ist danach nicht verloren, sondern wird sofort aufgelöst. Kennt jemand
den Effekt, wird die Person zur Komplizin gemacht (\enquote{Verraten Sie
es bitte noch nicht}). Misslingt der Force im Vorhersage-Trick, auf den
1089-Trick ausweichen; die Lektion bleibt dieselbe.
\end{vorbereitung}

\begin{ablauf}
\begin{description}
\item[0--10 · Staunmoment.] \enquote{Die unveränderliche Summe}: vier frei
  genannte Zahlen, die Summe trifft die versiegelte Vorhersage. Frage: Wo
  war hier die Freiheit?
\item[10--20 · Auflösung und Force-Begriff.] Die Rechnung an der Tafel;
  Einordnung als Force; Determinismus und Freiheits-Illusion
  (Newcomb-Paradoxon, Buch, Kap.~5b).
\item[20--45 · Theorie-Input.] Daten -- Befund -- Schlussfolgerung;
  Over-Interpretation, p-Hacking, HARKing, Garden of Forking Paths;
  Triangulation.
\item[45--60 · p-Hacking erleben.] Der Vorhersage-Trick (Buch,
  Abschnitt~10d); Tafel-Auflösung $1-(1-\alpha)^m$. Danach Übungsblatt
  Teil~B austeilen; die Tabelle in Aufgabe~2 vertieft die Formel (in der
  Stunde oder zu Hause).
\item[60--80 · Befund-Übung.] Rohdaten einer fiktiven Studie (Übungsblatt
  Teil~A, Aufgabe~1): deskriptive Befunde, zwei konkurrierende Interpretationen,
  entscheidende Zusatzdaten.
\item[80--90 · Selbstbezug.] Prüfung der eigenen bisherigen Ergebnisse:
  Habe ich nach Bestätigung gesucht? Was würde meine Hypothese
  herausfordern? (Übungsblatt Teil~B, Aufgabe~3)
\end{description}
\end{ablauf}

%==============================================================================
\section{Staunmoment: Die unveränderliche Summe}
%==============================================================================
\begin{staunbox}{Vier freie Zahlen}
Am Pult liegt seit Beginn der Stunde ein versiegelter Umschlag. Eine
Person nennt:
(1)~ihr Geburtsjahr, (2)~das Alter, das sie in diesem Jahr erreicht,
(3)~das Jahr eines Ereignisses, das ihr wichtig ist, (4)~wie viele
Silvester seitdem vergangen sind.

\emph{Jemand addiert. Dann öffnen wir den Umschlag.}
\end{staunbox}
\os{\newpage}

\begin{lehrhinweis}[Sprechtext (Vorschlag)]
\itshape
In diesem Umschlag liegt eine Zahl, die ich heute Morgen aufgeschrieben
habe. Ich brauche jemanden, der mir vier Zahlen über sich verrät, die ich
nicht kennen kann. Sagen Sie uns zuerst Ihr Geburtsjahr. \dots\ Und wie alt
werden Sie dieses Jahr an Ihrem Geburtstag -- oder sind es schon
geworden? \dots\ Jetzt denken Sie an ein Ereignis, das Ihnen wichtig war.
Sie müssen nicht sagen, was es war. Nur das Jahr. \dots\ Und wie viele
Silvester haben Sie seitdem gefeiert? \dots\ Das sind vier Zahlen, die nur
Sie kennen. Bitte addieren Sie sie (an eine zweite Person gewandt), und
Sie kontrollieren. \dots\ Lesen Sie die Summe laut vor. \dots\ Öffnen Sie
den Umschlag.
\end{lehrhinweis}

\begin{description}
\item[{\rot\bf Wo war hier die Freiheit?}]~\\[-\lblskip]
\bi
  \ite alle vier Zahlen waren frei gewählt
  \ite das Ergebnis stand trotzdem vorher fest
  \ite welche Erklärungen haben Sie?
  \ite was hätte die Vorhersage scheitern lassen können?
\ei
\os{\newpage}
\end{description}

%==============================================================================
\section{Auflösung: der Force}
%==============================================================================
\nsl{Die Auflösung ist eine Zeile Arithmetik. Geburtsjahr plus das im
laufenden Jahr erreichte Alter ergibt das laufende Jahr. Ereignisjahr plus
die seitdem vergangenen Jahre ergibt ebenfalls das laufende Jahr. Die
Summe ist also das Doppelte der Jahreszahl, gleichgültig, welche Person
und welches Ereignis gewählt wurden. In der Sprache der Zauberkunst ist
das ein \emph{Force}: Die Wahl ist frei, ihr Ergebnis ist es nicht. Das
Buch behandelt dasselbe Muster unter dem Stichwort Determinismus und
Freiheits-Illusion am Newcomb-Paradoxon \parencite{nozick1969newcomb}: Ein
Prädiktor füllt Boxen nach einer Vorhersage über eine Wahl, die erst
danach getroffen wird. Zum Zeitpunkt der Wahl liegt der Inhalt längst
fest.}
\begin{description}
\item[{\rot\bf Die Rechnung}]~\\[-\lblskip]
\[
\underbrace{\text{Geburtsjahr} + \text{Alter}}_{=\,J}
\;+\;
\underbrace{\text{Ereignisjahr} + \text{vergangene Jahre}}_{=\,J}
\;=\; 2J
\]
\bi
  \ite jede Eingabe war frei, die \emph{Struktur} der Rechnung nicht
  \ite die Freiheit hat das Ergebnis nie berührt
  \ite scheitern kann die Vorhersage nur an Messfehlern: falsches Alter,
    falsch gezählte Jahre
\ei
\os{\newpage}

\item[{\rot\bf Force}]~\\
\defn{force}{Ein \textbf{Force} ist ein Verfahren, das einer Person eine
freie Wahl lässt und trotzdem ein vorab festgelegtes Ergebnis
herbeiführt.}
\bi
  \ite Zauberkunst: Die versiegelte Vorhersage stimmt, weil das Verfahren
    sie erzwingt
  \ite Newcomb: Der Prädiktor hat entschieden, bevor die Person wählt
    \parencite{nozick1969newcomb}
  \ite Forschung: Eine Auswertung, die nur ein Ergebnis zulassen kann,
    prüft nichts
\ei
\os{\newpage}

\item[{\rot\bf Forces in der Forschung}]~\\[-\lblskip]
\bi
  \ite Stichprobe nur aus Freiwilligen, die das Thema mögen
  \ite Fragebogen nur mit zustimmend formulierten Items
  \ite \enquote{Erfolg} wird erst nach Sicht der Daten definiert
  \ite Auswertung so lange variieren, bis sie passt
  \ite Frage: Hätte meine Auswertung meine Hypothese überhaupt
    widerlegen können?
\ei
\nsl{Die letzte Frage verbindet die Einheit mit Horwichs Analyse des
Beweiswerts (LE~5 und LE~8): Nur Beobachtungen, die ein
Widerlegungsrisiko eingehen, haben vollen Beweiswert
\parencite{horwich1982probability}. Eine Auswertung, die strukturell
nur ein Ergebnis liefern kann, ist ein Force, kein Test.}
\os{\newpage}
\end{description}

%==============================================================================
\section{Daten, Befund, Schlussfolgerung}
%==============================================================================
\nsl{Der häufigste Fehler beim Interpretieren ist kein Rechenfehler, sondern
ein Ebenenwechsel. \emph{Daten} sind die erhobenen Werte. Ein
\emph{Befund} ist eine Aussage, die die Daten zusammenfasst und an ihnen
geprüft werden kann: Mittelwerte, Unterschiede, Zusammenhänge mit ihrer
Unsicherheit. Eine \emph{Schlussfolgerung} geht darüber hinaus; sie
erklärt den Befund, verallgemeinert ihn oder leitet eine Empfehlung ab.
Schlussfolgerungen sind notwendig, aber sie brauchen Annahmen, die man
offenlegen muss. Over-Interpretation entsteht, wenn eine Schlussfolgerung
wie ein Befund formuliert wird.}
\begin{description}
\item[{\rot\bf Drei Ebenen}]~\\[-\lblskip]
\begin{center}
\resizebox{0.95\textwidth}{!}{%
\begin{tikzpicture}[font=\sffamily,>=Stealth,
  st/.style={draw=rwuviolet,thick,fill=rwuviolet!8,rounded corners=3pt,
    minimum width=30mm,minimum height=16mm,align=center}]
  \node[st] (a) at (0,0) {\textbf{Daten}\\erhobene Werte};
  \node[st] (b) at (7,0) {\textbf{Befund}\\zusammengefasst,\\prüfbar};
  \node[st,draw=rwucyan,fill=rwucyan40!40] (c) at (14,0) {\textbf{Schlussfolgerung}\\Erklärung,\\Verallgemeinerung};
  \draw[->,very thick] (a) -- node[above]{Statistik} (b);
  \draw[->,very thick] (b) -- node[above]{Annahmen} (c);
\end{tikzpicture}}
\end{center}
\bi
  \ite Beispiel Monty Hall: 97 von 150 Wechselrunden gewonnen -- Wechseln
    gewann in 64,7\,\% -- Wechseln ist im Spiel überlegen
  \ite jede Stufe braucht etwas, das in der vorigen nicht steht
\ei
\nsl{Das Monty-Hall-Beispiel ist das konstruierte Rechenbeispiel aus LE~8.}
\os{\newpage}

\item[{\rot\bf Over-Interpretation}]~\\[-\lblskip]
\bi
  \ite Zusammenhang als Ursache lesen
  \ite über die Stichprobe hinaus verallgemeinern
  \ite \enquote{signifikant} als \enquote{groß} oder \enquote{wichtig} lesen
  \ite einen Einzelbefund als gesichert darstellen
  \ite aus dem Fehlen eines Befunds auf das Fehlen eines Effekts schließen
\ei
\os{\newpage}

\item[{\rot\bf Interpretation oder Spekulation?}]~\\[-\lblskip]
\bi
  \ite \textbf{Interpretation}: bleibt an den Befund gebunden, nennt
    ihre Annahmen, lässt Alternativen zu
  \ite \textbf{Spekulation}: geht über das hinaus, was die Daten tragen
    können -- erlaubt, wenn sie so heißt
  \ite Prüffrage: Welcher Befund würde meine Deutung widerlegen?
  \ite Pflicht: zu jeder Deutung mindestens eine Gegenhypothese
\ei
\os{\newpage}
\end{description}

%==============================================================================
\section{p-Hacking erleben: der Vorhersage-Trick}
%==============================================================================
\nsl{Für das p-Hacking hat das Buch eine Demonstration, in der die
Methode des Tricks und der zu vermittelnde Fehlschluss identisch sind
(Abschnitt~10d). Die Vorführende macht mehrere Vorhersagen über Zahlen,
die Freiwillige frei nennen. Beim Öffnen sind \emph{alle} Vorhersagen
falsch. Dann der Vorschlag, nicht auf die Treffer zu schauen, sondern auf
die Fehler: Die Abweichungen zwischen Vorhersage und Wirklichkeit ergeben
die Ziffern von $\pi$. Die Methode ist ein \emph{One-Ahead}: Die
Vorführende schreibt jede \enquote{Vorhersage für die nächste Runde} in
Wahrheit aus der gerade genannten Zahl plus der geplanten Abweichung;
die letzte Zahl ist geforct und die passende Karte vorab versiegelt. Das
Muster war also garantiert, weil jede Vorhersage nach dem Blick auf die
Daten konstruiert wurde.}
\begin{description}
\item[{\rot\bf Vier Vorhersagen, alle falsch}]~\\[-\lblskip]
\begin{center}
\begin{tabular}{@{}cccc@{}}
\toprule
\textbf{Runde} & \textbf{genannte Zahl $a_i$} & \textbf{Vorhersage $p_i$} & \textbf{Abweichung $p_i-a_i$} \\
\midrule
1 & 5 & 8 & \textbf{3} \\
2 & 9 & 10 & \textbf{1} \\
3 & 2 & 6 & \textbf{4} \\
4 & 7 & 8 & \textbf{1} \\
\bottomrule
\end{tabular}
\end{center}
\bi
  \ite keine Vorhersage trifft -- aber die Fehler lauten 3, 1, 4, 1
\ei
\nsl{Die Tabelle ist das durchgerechnete Beispiel aus dem Buch; in der
Stunde stehen dort die tatsächlich genannten Zahlen.}
\os{\newpage}

\item[{\rot\bf Die Methode ist der Fehlschluss}]~\\[-\lblskip]
\bi
  \ite jede Vorhersage wurde \emph{nach} der genannten Zahl geschrieben:
    $p_i = a_i + r_i$
  \ite das Muster war garantiert, egal welche Zahlen kamen
  \ite genau das ist \textbf{p-Hacking}: Bedeutung wird nachträglich in
    Daten hineingerechnet
  \ite dafür braucht es keinen bösen Willen, nur Auswahl nach Sicht der
    Daten
\ei
\os{\newpage}
\end{description}

%==============================================================================
\section{Viele Wege, ein Treffer}
%==============================================================================
\nsl{In der Forschung braucht es für p-Hacking keinen Trick. Es genügt,
nach Sicht der Daten unter vielen plausiblen Auswertungen die eine zu
wählen, die \enquote{funktioniert}. \textcite{gelmanloken2014forking}
nennen das den \emph{Garden of Forking Paths}: Jede Entscheidung über
Ausreißer, Teilgruppen, Kovariaten, Messzeitpunkte oder Abbruchkriterien
ist eine Weggabelung, und wer erst am Ende des Gartens entscheidet,
welchen Weg er gegangen sein will, findet fast immer einen Treffer.
\textcite{simmons2011falsepositive} haben gezeigt, dass solche
\enquote{researcher degrees of freedom} die Falsch-Positiv-Rate weit über
die nominellen 5\,\% treiben. \textcite{kerr1998harking} hat den
verwandten Fall benannt, dass eine Hypothese erst nach den Ergebnissen
gebildet und dann als vorab aufgestellt präsentiert wird (HARKing). Und
\textcite{ioannidis2005false} hat argumentiert, dass bei kleinen
Studien, kleinen Effekten und vielen getesteten Beziehungen ein großer
Teil der veröffentlichten positiven Befunde falsch sein dürfte.}
\begin{description}
\item[{\rot\bf Begriffe}]~\\[-\lblskip]
\bi
  \ite \textbf{p-Hacking}: Auswertung variieren, bis $p<0{,}05$
    \parencite{simmons2011falsepositive}
  \ite \textbf{HARKing}: Hypothese nach den Ergebnissen bilden und als
    vorab ausgeben \parencite{kerr1998harking}
  \ite \textbf{Garden of Forking Paths}: viele Gabelungen, auch ohne
    bewusstes Probieren \parencite{gelmanloken2014forking}
  \ite \textbf{Data Dredging}: in großen Datensätzen ohne Hypothese nach
    Mustern graben
\ei
\os{\newpage}

\item[{\rot\bf Die Tafel-Auflösung}]~\\[-\lblskip]
\[
P(\text{mindestens ein Treffer}) \;=\; 1-(1-\alpha)^m
\]
\begin{center}
\small
\begin{tabular}{@{}lcccccccc@{}}
\toprule
$m$ & 1 & 2 & 5 & 10 & 14 & 20 & 50 & 100 \\
\midrule
$\alpha=0{,}05$ & 5\,\% & 9,8\,\% & 22,6\,\% & 40,1\,\% & 51,2\,\% & 64,2\,\% & 92,3\,\% & 99,4\,\% \\
\bottomrule
\end{tabular}
\end{center}
\bi
  \ite $m$: Zahl der Versuche (Auswertungen, Hypothesen,
    Zwischenstände)
  \ite ab $m=14$ ist ein Zufallstreffer wahrscheinlicher als keiner
  \ite $m$ ist die \enquote{Magier-Freiheit} des Tricks
\ei
\nsl{Die Werte sind nachgerechnet; der Wert für $m=14$ steht so auch im
Buch. Eine strengere Schwelle wie $p<0{,}005$ verschiebt das Problem nur
\parencite{benjamin2018redefine}: Bei $m=100$ liegt die Wahrscheinlichkeit
für mindestens einen Zufallstreffer dann immer noch bei 39,4\,\%.}
\os{\newpage}

\item[{\rot\bf Die Bayes-Pointe}]~\\[-\lblskip]
\bi
  \ite viele Hypothesen teilen sich den Vorab-Glauben, jede bekommt nur
    wenig
  \ite ein Treffer unter vielen war ohnehin zu erwarten
  \ite also bleibt die Posteriori nahe der Priori
  \ite 100 versiegelte Vorhersagen, eine trifft: Erwartung, kein Wunder
\ei
\nsl{Ein Verfahren, das \emph{immer} ein hübsches Ergebnis liefert, ist
so wenig Evidenz wie ein Test, der die Nullhypothese stets ablehnt (Buch,
Abschnitt~10d; \textcite{ioannidis2005false}).}
\os{\newpage}

\item[{\rot\bf Ein Lehrstück aus dem Buch}]~\\[-\lblskip]
\bi
  \ite Monty-Hall-Feldstudie, Teilnehmende \emph{ohne} Vorwissen
  \ite Trefferquote U1: 25\,\%, U2: 60\,\%, U3: 83\,\%
  \ite Differenz U3 gegen U1: $+58$ Prozentpunkte $[+6;\,+81]$
  \ite aber: sechs gegen acht Personen, Analyse erst nach Sichtung der
    Daten gerechnet
\ei
\nsl{Das Buch kommentiert diesen Befund selbst als den, dem man am
wenigsten trauen darf, und ergänzt: \enquote{Genau so entstehen
Ergebnisse, die in der Replikation verschwinden} (Kap.~12f). Die
Schichtung nach Vorwissen war nicht vorab geplant und ist im Buch
ausdrücklich als explorativ markiert. Das ist ein Beispiel für saubere
Praxis: Der Befund wird berichtet, aber als Hypothese für eine
Folgestudie, nicht als Ergebnis.}
\os{\newpage}
\end{description}

%==============================================================================
\section{Gegenmittel}
%==============================================================================
\nsl{Gegen die Freiheiten der Auswertung helfen drei Dinge. Erstens die
\emph{Präregistrierung}: Hypothesen, Stichprobengröße und Auswertungsplan
werden vor der Datenerhebung öffentlich festgelegt. Damit wird sichtbar,
was vorhergesagt und was nachträglich gefunden wurde; beides ist
legitim, solange man es unterscheidet \parencite{nosek2018preregistration}.
Die Präregistrierung ist das Gegenteil des One-Ahead: Die Vorhersage
liegt tatsächlich vorher im Umschlag. Zweitens die \emph{Triangulation}:
Derselbe Gegenstand wird mit verschiedenen Daten, Methoden oder Personen
untersucht \parencite{denzin1970research}. Ein Befund, der nur in einer
Auswertung auftaucht, ist verdächtig; einer, den verschiedene Zugänge
unabhängig stützen, ist robuster. Drittens die \emph{Gegenhypothese}: Wer
zu jeder Deutung aktiv nach der besten Alternative sucht, schützt sich
gegen den Bestätigungsfehler (LE~7).}
\begin{description}
\item[{\rot\bf Präregistrierung}]~\\[-\lblskip]
\bi
  \ite \emph{vor} der Erhebung festlegen: Hypothesen, $n$, Ausschlüsse,
    Auswertung
  \ite öffentlich und mit Zeitstempel hinterlegen
  \ite Abweichungen sind erlaubt, müssen aber berichtet werden
  \ite trennt Vorhersage von Nachhersage \parencite{nosek2018preregistration}
  \ite das Gegenteil des One-Ahead: Die Vorhersage liegt wirklich vorher
    im Umschlag
\ei
\os{\newpage}

\item[{\rot\bf Triangulation}]~\\[-\lblskip]
\bi
  \ite \textbf{Daten}: verschiedene Quellen, Zeitpunkte, Orte
  \ite \textbf{Forschende}: mehrere Personen werten unabhängig aus (LE~9)
  \ite \textbf{Methoden}: Zahlen und Interviews, Experiment und Beobachtung
  \ite \textbf{Theorien}: denselben Befund aus mehreren Perspektiven
    deuten \parencite{denzin1970research}
\ei
\os{\newpage}

\item[{\rot\bf Prüffragen an die eigene Auswertung}]~\\[-\lblskip]
\bi
  \ite Stand meine Auswertung fest, bevor ich die Daten sah?
  \ite Wie viele Wege habe ich ausprobiert, wie viele berichte ich?
  \ite Was ist die beste Gegenhypothese?
  \ite Welche Zusatzdaten würden zwischen beiden entscheiden?
  \ite Hätte meine Auswertung meine Hypothese widerlegen können?
\ei
\os{\newpage}
\end{description}

\begin{merkbox}
Ein Muster, das jede Datenlage hervorbringen würde, sagt nichts über die
konkrete Datenlage. Wer erst nach dem Blick auf die Daten entscheidet,
wonach er sucht, hat keinen Befund, sondern eine Hypothese für die
nächste Studie.
\end{merkbox}
\os{\newpage}

%==============================================================================
\section{Übung in der Stunde}
%==============================================================================
\begin{description}
\item[{\rot\bf Befund-Übung (Kleingruppen, 20 Min.)}]~\\[-\lblskip]
\bi
  \ite fiktive Studie: Tagebuch des Staunens und Klausurpunkte
    (Übungsblatt Teil~A)
  \ite (a) Formulieren Sie die deskriptiven Befunde
  \ite (b) Formulieren Sie zwei konkurrierende Interpretationen
  \ite (c) Welche Zusatzdaten würden zwischen ihnen entscheiden?
  \ite (d) Wie viele Wege hätte man auswerten können?
\ei
\os{\newpage}

\item[{\rot\bf Selbstbezug (Einzelarbeit, 10 Min.)}]~\\[-\lblskip]
\bi
  \ite Notieren Sie Ihr bisher wichtigstes Ergebnis im Semesterprojekt
  \ite Haben Sie gezielt nach Bestätigung gesucht?
  \ite Welcher Befund würde Ihre Hypothese herausfordern?
  \ite Wo liegen Ihre eigenen Weggabelungen?
\ei
\os{\newpage}
\end{description}

\begin{lehrhinweis}
Die nachgerechneten Kennwerte der fiktiven Studie stehen im
Erwartungshorizont am Ende von Übungsblatt Teil~B. Der Datensatz ist so konstruiert,
dass der Unterschied zwischen den Tagebuch-Gruppen groß und
\enquote{hochsignifikant} ist, Tagebuch und Anwesenheit aber so eng
zusammenhängen, dass sich ihre Wirkungen mit diesen Daten nicht trennen
lassen. Die zentrale Einsicht: Die konkurrierende Interpretation
(Engagement als Drittvariable) ist mit den Daten genauso verträglich, und
nur ein anderes Design, etwa die zufällige Zuteilung zum Tagebuch,
könnte entscheiden. Im Selbstbezug genügt es, wenn jede Person eine
eigene Weggabelung konkret benennt.
\end{lehrhinweis}

%==============================================================================
\section{Lernziele und Aufgaben}
%==============================================================================
\begin{description}
\item[{\rot\bf Nach dieser Einheit können Sie}]~\\[-\lblskip]
\bi
  \ite Daten, Befund und Schlussfolgerung unterscheiden
  \ite zu einem Befund konkurrierende Interpretationen formulieren und
    entscheidende Zusatzdaten benennen
  \ite p-Hacking, HARKing und den Garden of Forking Paths erklären
  \ite die Wahrscheinlichkeit von Zufallstreffern bei $m$ Tests berechnen
  \ite Präregistrierung und Triangulation als Gegenmittel einsetzen
\ei
\os{\newpage}

\item[{\rot\bf Bis zur nächsten Woche}]~\\[-\lblskip]
\bi
  \ite \textbf{Befund-Übung} fertigstellen (Übungsblatt Teil~A, Aufgabe~1)
  \ite \textbf{Reflexion} zum Semesterprojekt: Haben Sie Ergebnisse
    gesucht, die Ihre Hypothese \emph{bestätigen}? Welche Befunde könnten
    sie \emph{herausfordern}? (Teil~B, Aufgabe~3)
  \ite \textbf{Tagebuch des Staunens}: eine Entscheidung, die sich frei
    anfühlte und es vielleicht nicht war
  \ite für LE~11: einen Absatz aus Ihrem Projekt mitbringen, in dem Sie
    ein Ergebnis deuten
\ei
\os{\newpage}
\end{description}

\handoutinput{handouts/LE10}
